\documentclass[11pt]{article}
\usepackage[a4paper,margin=27mm]{geometry}
\usepackage[T1]{fontenc}
\usepackage{lmodern,microtype,amsmath,amssymb,amsthm,mathtools,enumitem}
\usepackage[hidelinks,hypertexnames=false]{hyperref}
\newtheorem{theorem}{Theorem}
\newtheorem{lemma}[theorem]{Lemma}
\newtheorem{proposition}[theorem]{Proposition}
\newtheorem{corollary}[theorem]{Corollary}
\theoremstyle{remark}\newtheorem{remark}[theorem]{Remark}
\newcommand{\R}{\mathbb R}
\newcommand{\Dscr}{\mathcal D}
\newcommand{\Escr}{\mathcal E}
\newcommand{\Ascr}{\mathcal A}
\newcommand{\supp}{\operatorname{supp}}
\newcommand{\one}{\mathbf 1}
\title{What Closed Competitive Choice Requires of Independent Rankings}
\author{}\date{9 October 2026}
\begin{document}\maketitle

\noindent
This note continues \emph{Independent Rankings of Binary Advice} (v3) and
\emph{Beyond Additive Guidance Rankings} (BA). It does not construct a ranking
satisfying (S) together with the other requirements for $n\ge3$ worlds. That
problem remains open. It proves four things that narrow it.
\begin{enumerate}[leftmargin=*,itemsep=2pt]
\item The additive obstruction needs an additive weak extension only for
comparisons between \emph{regular} models: models that never contradict
themselves and are correct on unanimous problems
(Theorem~\ref{thm:regular}). Comparisons involving the selected models are
not the right class: an explicit ranking satisfies (S) although those
comparisons have no additive weak extension (Proposition~\ref{prop:priority}).
The non-additive witness in BA uses non-regular models, and it collapses on
regular models (Remark~\ref{rem:collapse}).
\item Continuity forces additivity. If, on regular models, each world's ranking
is complete and continuous in measure, or an intersection of such rankings,
then it has an additive weak extension. So the vacuous model is dominated
(Theorem~\ref{thm:continuous}).
\item (S) forces specific structure: face support of competitive coherent
models, uncountably many distinct selected models on every edge, and limits
on real-valued primary criteria (Section~\ref{sec:structure}).
\item A remaining solution must be discontinuous or incomplete in a specific
sense. We conjecture that none exists (Section~\ref{sec:remaining}).
\end{enumerate}

\section{Setting}

We use the setting of v3 and BA without change. There are $n$ worlds and
$\Delta=\{p\ge0:\sum_ip_i=1\}$. Models are measurable acceptance sets
$\Dscr\subseteq\R^n$, identified modulo null sets. Binary problems are pairs
$\{g,-g\}$ represented by $g\in H=\{g:g_1>0\}$, and $\mu$ is standard Gaussian
measure. The advice $a_\Dscr(g)\in\{+,0,-,\times\}$ has guidance rank
$r_i(g,\Dscr)\in\{0,1,2,3\}$ at world $i$ (good, silent, bad, contradictory).
We write $\Dscr\unlhd_i\Escr$ when $r_i(\cdot,\Dscr)\le r_i(\cdot,\Escr)$ a.e.
A ranking is a family of preorders $\preceq_i$, smaller meaning better.
The requirements (E), (G), (IB), (LA), (A), (S), (C) are those of v3 and BA.
An \emph{additive guidance loss} is
$L_i(\Dscr)=\int_H\ell_i(g,a_\Dscr(g))\,d\mu$ with
$0=\ell_i(\mathrm{good})<\ell_i(0)<\ell_i(\mathrm{bad})<\ell_i(\times)$ a.e.
and an integrable envelope.

A problem is \emph{unanimous} if all its payoffs have the same sign, and
\emph{mixed} otherwise. Write $H'$ for the mixed problems in $H$.

A model is \emph{regular} if, up to null sets, it is correct on every
unanimous problem and never gives contradictory advice. Every coherent model
is regular: it accepts every nonzero nonnegative gamble, and it cannot accept
both $g$ and $-g$, since their sum is $0$.

The \emph{regularization} $\Escr^\circ$ of a model $\Escr$ gives the correct
advice on unanimous problems, silence where $\Escr$ contradicts itself on a
mixed problem, and the advice of $\Escr$ elsewhere. It is regular and
satisfies $\Escr^\circ\unlhd_i\Escr$ at every world.

We use BA's Lemma~1 (from (LA) and (S) by the simplex lemma):
\begin{enumerate}[label=(KKM\alph*),leftmargin=*]
\item every model is weakly dominated at all worlds by some selected $C_q$;
\item every undominated model is tied at every world with some selected
$C_p$, and is competitive on $\supp p$.
\end{enumerate}

\section{Regular comparisons suffice for the obstruction}

\begin{theorem}\label{thm:regular}
Let $n\ge3$. Suppose a ranking of all measurable models satisfies (G), (LA)
and (S). Suppose also that for each world $i$ some additive guidance loss
$L_i$ satisfies
\begin{equation}
 \Dscr\preceq_i\Escr\ \Longrightarrow\ L_i(\Dscr)\le L_i(\Escr)
 \qquad\text{for all regular }\Dscr,\Escr.
 \label{eq:regular-extension}
\end{equation}
Then every $\Dscr_M$ whose credal set contains three simplex vertices is
dominated by a coherent model. In particular the vacuous model is dominated.
\end{theorem}
\begin{proof}
\emph{Coherent Bayes models at every prior.} Fix $p\in\Delta$ and choose, at
each problem, an advice minimizing $\sum_ip_i\ell_i(g,\cdot)$, breaking ties
in a fixed order. Call the resulting model $B$. It is $p$-Bayes for $L$.

$B$ is regular. On a unanimous problem, correct advice costs $0$ at every
world and any other advice costs more. Contradictory advice costs strictly
more than silence at every world.

By (KKMa), some selected $C_q$ has $C_q\preceq_jB$ at every world. Both
models are regular, so \eqref{eq:regular-extension} gives
$L_j(C_q)\le L_j(B)$ for all $j$. Hence $C_q$ is a coherent $p$-Bayes model.

\emph{The target is Bayes somewhere.} Suppose $\Dscr_M$ is undominated. By
(KKMb) it is competitive on $\supp p$ for some $p$. Take any model $\Escr$.
Competitiveness gives some $j$ with $\Dscr_M\preceq_j\Escr^\circ$. Both models
are regular, and penalties are ordered by guidance rank. So
\[
 L_j(\Dscr_M)\le L_j(\Escr^\circ)\le L_j(\Escr).
\]
Thus no model is strictly better than $\Dscr_M$ for $L$ at every world. By
v3's Lemma~11(b), which holds for every additive guidance loss,
$\Dscr_M$ is $L$-Bayes at some prior.

\emph{Contradiction.} Coherent $L$-Bayes models exist at every prior, so v3's
Proposition~13 forbids this. Hence some $\Escr$
dominates $\Dscr_M$. By (KKMa), some selected $C_q$ satisfies
$C_q\preceq_j\Escr$ at every world, so $C_q$ dominates $\Dscr_M$.
\end{proof}

Theorem~\ref{thm:regular} strengthens BA's Theorem~2, which assumed
\eqref{eq:regular-extension} for all models. The regularization step also
supplies the line ``hence competitive for $L$'' in BA's proof, which needs an
argument when $L$ respects only some comparisons.

\begin{remark}[The six-region witness is not regular]\label{rem:collapse}
In BA's six-region construction, each region $B_j$ contains unanimous and
mixed problems. Its non-additivity witnesses are silent on whole regions, so
they are not regular. On regular models, the corresponding family compares by
inclusion.

To see this, fix a world $i$ and $A\subseteq\{1,\ldots,6\}$. Let $F_A$ be regular,
silent on the mixed problems of $B_j$ for $j\in A$, and correct at $i$ on the
other mixed problems. At a sampled tuple, let $S$ be the set of coordinates
falling in mixed problems. The rank vector of $F_A$ at world $i$ is then
$\one_{A\cap S}$. Every $S$ occurs with positive probability, since the
draws are independent and each region contains both kinds of problem. So
$F_A\preceq_iF_B$ requires $A\cap S\le_QB\cap S$ for every $S$.

For $j\in A\setminus B$, take $S=\{j\}$. This would need $\{j\}\le_Q\varnothing$,
which is false. Conversely, inclusion implies the comparison. So the
certificate in BA (four comparisons summing to $-e_1$) has no regular
counterpart in this family.

Restricting the regions to mixed problems restores the witness on regular
models: the rank vector is $\one_A$ at every tuple. But BA's proof that
coherent models are undominated pairs a region with unanimous problems in
the other regions, and so does not apply as stated.
\end{remark}

The obvious alternative to \eqref{eq:regular-extension} uses comparisons
involving the selected models. That class does not work: an additive weak
extension of those comparisons can fail to exist while (S) holds. BA's final
section suggested this route.

\begin{proposition}\label{prop:priority}
Let $L$ be an additive guidance ranking satisfying the hypotheses of v3's
Theorem~4, with selection $C_p$. Write $U$ for the unanimous problems
in $H$, and put
\[
 P_i(\Dscr)\coloneqq\mu\{g\in U:r_i(g,\Dscr)\ge1\},
\]
\[
 \Dscr\preceq'_i\Escr\iff P_i(\Dscr)<P_i(\Escr)\ \text{ or }\
 \bigl(P_i(\Dscr)=P_i(\Escr)\text{ and }L_i(\Dscr)\le L_i(\Escr)\bigr).
\]
Then $\preceq'$ satisfies (G), (IB), (LA), (S) and (C) with the same
selection. Moreover, no additive guidance loss $\Lambda$ satisfies
$\Lambda_i(C_p)\le\Lambda_i(\Escr)$ whenever $C_p\preceq'_i\Escr$.
\end{proposition}
\begin{proof}
\emph{Preorder, (IB), (G).} A lexicographic combination of two complete
preorders is a complete preorder. Both $P_i$ and $L_i$ are integrals of local
penalties, so comparisons depend only on the problems where two models
differ; this is (IB).

A weak guidance improvement weakly decreases both $P_i$ and $L_i$. A strict
one strictly decreases $L_i$ and weakly decreases $P_i$. Either way the
improvement is strict for $\preceq'_i$.

\emph{(S).} Every coherent model is correct on $U$, so $P_i(C_p)=0$. Fix
$\Escr$. If $P_i(\Escr)>0$ for some $i\in\supp p$, then $C_p\prec'_i\Escr$.
Otherwise the comparisons at worlds in $\supp p$ are those of $L$, and
competitiveness for $L$ supplies some $i\in\supp p$ with
$L_i(C_p)\le L_i(\Escr)$. Closedness holds because
\[
 \{p:C_p\preceq'_i\Escr\}=
 \begin{cases}
 \Delta, & P_i(\Escr)>0,\\
 \{p:L_i(C_p)\le L_i(\Escr)\}, & P_i(\Escr)=0.
 \end{cases}
\]
The second set is closed by v3's Theorem~4(b).

\emph{(LA), (C).} A model weakly dominating $C_p$ has $P\equiv0$, so it
weakly dominates $C_p$ for $L$. Hence $C_p$ is undominated, and a model
tied with $C_p$ is $L$-tied with it and therefore agrees with $C_p$ a.e.

\emph{No additive weak extension.} Suppose $\Lambda$ were one. Fix an
interior $p$. Since $C_p$ cannot agree a.e.\ with both ideal models $D_1$ and
$D_2$, choose $i\in\{1,2\}$ with $C_p\not\approx D_i$. Let $\Escr_\varepsilon$ be $D_i$ made
silent on a set $Z_\varepsilon\subseteq U$ with $\mu(Z_\varepsilon)=\varepsilon$.
Then $P_i(\Escr_\varepsilon)>0$, so $C_p\prec'_i\Escr_\varepsilon$, and
\[
 \Lambda_i(C_p)\le\Lambda_i(\Escr_\varepsilon)
 \le\Lambda_i(D_i)+\int_{Z_\varepsilon}h\,d\mu,
\]
where $h$ is the integrable envelope. Letting $\varepsilon\to0$ gives
$\Lambda_i(C_p)\le\Lambda_i(D_i)$. But $D_i$ is a strict guidance
improvement of $C_p$ at $i$, so strictly ordered penalties give
$\Lambda_i(D_i)<\Lambda_i(C_p)$. This is a contradiction.
\end{proof}

Proposition~\ref{prop:priority} covers exactly the targets its additive base
covers, so it does not bear on full coverage. Its role is diagnostic: it
satisfies (S) although its comparisons involving selected models have no
additive weak extension. On regular models it coincides with $L$, as
Theorem~\ref{thm:regular} requires.

\section{Continuity forces additivity}

Fix a world $i$. A regular model is determined by its rank function at $i$
on mixed problems: rank $0$ means the advice correct at $i$, rank $1$ silence,
rank $2$ the advice incorrect at $i$. Every measurable rank function arises.
So the regular models correspond bijectively to
\[
 \Ascr\coloneqq\{x:H'\to\{0,1,2\}\text{ measurable}\}\ \text{modulo null sets},
 \qquad \delta(x,y)\coloneqq\mu\{x\ne y\}.
\]
For $Q\subseteq H'$ write $x\one_Q$ for the function equal to $x$ on $Q$ and
$0$ elsewhere. A preorder $\sqsubseteq$ on $\Ascr$ is
\begin{itemize}[leftmargin=*,itemsep=2pt]
\item \emph{independent} if, whenever $x=x'$ and $y=y'$ on $Q$ while $x=y$ and
$x'=y'$ off $Q$, we have $x\sqsubseteq y\iff x'\sqsubseteq y'$;
\item \emph{weakly monotone} if $x\le y$ a.e.\ implies $x\sqsubseteq y$, and
\emph{strictly monotone} if moreover $x\sqsubset y$ when $x\ne y$ on a set of
positive measure;
\item \emph{continuous} if $\{x:x\sqsubset y\}$ and $\{x:y\sqsubset x\}$ are
$\delta$-open for every $y$.
\end{itemize}
These are (IB), weak or strict (G), and continuity in measure, transcribed
for regular models.

\begin{lemma}\label{lem:debreu}
Let $\sqsubseteq$ be a complete, independent, weakly monotone, continuous
preorder on $\Ascr$. Then there are measurable $0\le\ell_1\le\ell_2$ on $H'$
with $\int\ell_2\,d\mu<\infty$ such that, with $\ell_0\coloneqq0$,
\[
 x\sqsubseteq y\iff\int_{H'}\ell_{x(g)}(g)\,d\mu\le\int_{H'}\ell_{y(g)}(g)\,d\mu .
\]
If $\sqsubseteq$ is strictly monotone, then $0<\ell_1<\ell_2$ a.e.
\end{lemma}
\begin{proof}
\emph{Step 1: inessential sets.} Call $Q$ \emph{inessential} if $x\sim y$
whenever $x=y$ off $Q$.
\begin{itemize}[leftmargin=*,itemsep=2pt]
\item A finite union of inessential sets is inessential. If $x=y$ off
$Q_1\cup Q_2$, let $z=y$ on $Q_1$ and $z=x$ elsewhere. Then $x\sim z$, since
they differ only on $Q_1$, and $z\sim y$, since they differ only on
$Q_2\setminus Q_1$.
\item A countable union of inessential sets is inessential. If $x=y$ off
$\bigcup_mQ_m$, let $z_m=y$ on $Q_1\cup\dots\cup Q_m$ and $z_m=x$ elsewhere. Then
$x\sim z_m$ and $\delta(z_m,y)\to0$. If $x\sqsubset y$, the open set
$\{w:x\sqsubset w\}$ contains $y$, hence some $z_m$, which is impossible.
The case $y\sqsubset x$ is symmetric. By completeness, $x\sim y$.
\end{itemize}
Choose inessential sets whose measures approach the supremum of all such
measures, and let $N_0$ be their union. Then $N_0$ is inessential, and every
non-null subset of $S\coloneqq H'\setminus N_0$ is essential: otherwise its
union with $N_0$ would exceed the supremum.

Every $x$ satisfies $x\sim x\one_S$, since these differ only on $N_0$. So
$\sqsubseteq$ is determined by its restriction to functions on $S$. If
$\mu(S)=0$, take $\ell_1=\ell_2=0$. Under strict monotonicity, $N_0$ is null:
otherwise $0\sqsubset\one_{N_0}$ would contradict inessentiality.

\emph{Step 2: Debreu's theorem.} Let $P$ be a finite partition of $S$ into
$k\ge3$ non-null pieces. Then $\Ascr_S$ is the product of the spaces
$\Ascr_Q$, $Q\in P$, and $\delta$ is the sum of the piece metrics, so it
induces the product topology.

Each $\Ascr_Q$ is separable: simple functions over a countable generating
algebra are dense. Each is path-connected: given $x,y$, let
$x_t=y$ on $Q\cap\{\Phi(g_2)<t\}$ and $x_t=x$ elsewhere on $Q$, where $\Phi$ is
the standard normal distribution function. Then $t\mapsto x_t$ is continuous
from $x_0=x$ to $x_1=y$.

Every piece is essential, and independence holds for every union of pieces.
So $\sqsubseteq$ is a continuous weak order on a product of at least three
connected separable essential factors, with all subsets of factors
independent. Debreu's theorem \cite{Debreu,Wakker} gives continuous
$u_Q$ with
\[
 x\sqsubseteq y\iff\sum_{Q\in P}u_Q(x|_Q)\le\sum_{Q\in P}u_Q(y|_Q).
\]
This representation is unique up to a common positive scale and additive
constants.

Fix one such partition $P_0$ with representation $U$, and let $P$ be any finite
partition of $S$. Let $P'$ be the common refinement of $P$ and $P_0$, with
null cells discarded. Debreu's theorem for $P'$ gives a continuous
$P'$-additive representation $W$. Grouping its terms along the cells of $P_0$
makes $W$ a continuous $P_0$-additive representation, so $W=aU+b$ with $a>0$.
Hence $U$ is $P'$-additive, and grouping along $P$ shows it is $P$-additive.
Normalizing,
\begin{equation}
 U(x)-U(0)=\sum_{Q\in P}\bigl(U(x\one_Q)-U(0)\bigr)
 \qquad\text{for every finite partition }P\text{ of }S.
 \label{eq:partition-additive}
\end{equation}

\emph{Step 3: densities.} For measurable $Q\subseteq S$ and $k=1,2$, put
$\nu_k(Q)\coloneqq U(k\one_Q)-U(0)$.
\begin{itemize}[leftmargin=*,itemsep=2pt]
\item Applying \eqref{eq:partition-additive} to $k\one_{Q\cup Q'}$ with the
partition $\{Q,Q',S\setminus(Q\cup Q')\}$ shows that $\nu_k$ is finitely
additive.
\item If $Q_m\downarrow$ with null intersection, then
$\delta(k\one_{Q_m},0)\to0$, and continuity of $U$ gives $\nu_k(Q_m)\to0$.
So $\nu_k$ is countably additive.
\item Weak monotonicity gives $0\le\nu_1\le\nu_2$, and both vanish on null sets.
\end{itemize}
The Radon--Nikodym theorem gives densities $\ell_k$ on $S$ with
$0\le\ell_1\le\ell_2$ a.e.\ and $\int\ell_2=\nu_2(S)<\infty$. Set $\ell_k=0$ on
$N_0$. For any $x$, apply \eqref{eq:partition-additive} to the partition of
$S$ by the value of $x$:
\[
 U(x\one_S)-U(0)=\nu_1\{x=1\}+\nu_2\{x=2\}=\int\ell_{x(g)}(g)\,d\mu .
\]
Since $x\sim x\one_S$, this proves the representation.

Under strict monotonicity, $0\sqsubset\one_Q$ and $\one_Q\sqsubset2\one_Q$ for
every non-null $Q$. Hence $\int_Q\ell_1>0$ and $\int_Q(\ell_2-\ell_1)>0$ for
every such $Q$, which gives $0<\ell_1<\ell_2$ a.e.
\end{proof}

\begin{theorem}\label{thm:continuous}
Let $n\ge3$, and let a ranking of all measurable models satisfy (G), (LA)
and (S). Suppose that, at each world $i$, the restriction of $\preceq_i$ to
regular models is the intersection of a family $\{\sqsubseteq_i^\alpha\}$ of
complete, independent, weakly monotone, continuous preorders. Then every
$\Dscr_M$ whose credal set contains three simplex vertices is dominated by a
coherent model. This applies in particular when each $\preceq_i$ is complete
and continuous on regular models.
\end{theorem}
\begin{proof}
Fix $i$. Lemma~\ref{lem:debreu} gives densities $\ell^\alpha_1\le\ell^\alpha_2$
for each $\alpha$; put $\ell^\alpha_0=0$.

For $k=1,2$, the sets $\{\ell^\alpha_k>\ell^\alpha_{k-1}\}$ cover $H'$ up to a
null set. Otherwise some non-null $Q$ misses all of them. The regular models
with rank functions $(k-1)\one_Q$ and $k\one_Q$ differ by a strict guidance
improvement, so by strict (G) the second is not ranked at least as well as
the first. Hence some $\sqsubseteq^\alpha_i$ ranks it strictly worse, giving
$\int_Q(\ell^\alpha_k-\ell^\alpha_{k-1})>0$. This is a contradiction.

A family of measurable sets has a countable subfamily with the same essential
union (BA, proof of the Corollary). So there are $\alpha_1,\alpha_2,\ldots$ for
which
\[
 \ell_k\coloneqq\sum_mc_m\ell^{\alpha_m}_k,\qquad
 c_m\coloneqq\frac{2^{-m}}{1+\int\ell^{\alpha_m}_2\,d\mu},
\]
satisfy $0<\ell_1<\ell_2$ a.e.\ and $\int\ell_2\le1$.

Define local penalties at world $i$. On mixed problems, ranks $0,1,2,3$ cost
$0,\ell_1,\ell_2,\ell_2+1$. On unanimous problems they cost $0,1,2,3$. This is
an additive guidance loss $L_i$.

For regular $\Dscr\preceq_i\Escr$, every $\sqsubseteq^{\alpha_m}_i$ compares
them the same way, so each integral is weakly ordered, and therefore so is
their weighted sum. The unanimous contributions are both zero. Thus
\eqref{eq:regular-extension} holds, and Theorem~\ref{thm:regular} applies.
\end{proof}

Every additive guidance loss is continuous in measure:
$|L_i(\Dscr)-L_i(\Escr)|\le\int_Zh_i\,d\mu$, where $Z$ is the set of problems on
which the models differ. So Theorem~\ref{thm:continuous} contains the additive
case and BA's Corollary. Its content is that no complete continuous ranking
escapes, however non-additive its definition looks: Lemma~\ref{lem:debreu}
makes it additive on regular models.

\begin{remark}[A closed graph is not enough]\label{rem:closed-graph}
Theorem~\ref{thm:continuous} needs open strict contour sets. A weaker notion
of continuity does not suffice. Take BA's six-region tuple ranking, but with
the six regions partitioning $H'$, and intersect it with the background
comparison $\unlhd_i$ restricted to unanimous problems.
\begin{itemize}[leftmargin=*,itemsep=2pt]
\item Its graph is closed in measure. If models converge in measure, their
rank vectors converge in product measure, hence a.e.\ along a subsequence.
Ranks take finitely many values, so a.e.\ inequalities pass to the limit.
\item It satisfies (IB) and strict (G). For the tuple part this is BA's
argument; the unanimous part is a pointwise comparison. A strict improvement
is strict in at least one of the two parts and weak in the other.
\item On regular models the unanimous parts agree, so it reduces to the tuple
ranking.
\item It is not an intersection as in Theorem~\ref{thm:continuous}. The regular
models $F_A$ of Remark~\ref{rem:collapse} now have rank vector $\one_A$ at
every tuple, so they compare exactly as $Q$. The proof of
Theorem~\ref{thm:continuous} would produce strictly positive weights
$w_j=\int_{B_j}\ell_1$ respecting $Q$, contradicting BA's certificate.
\end{itemize}
This ranking fails (S) by BA's finite-sample theorem.
\end{remark}

\section{Structure forced by (S)}\label{sec:structure}

In the additive case, (S) produces Bayes models at every prior. Without
additivity no prior is available. Competitiveness of $C_p$ depends only on
$\supp p$, so the numerical index enters (S) only through closedness. This
section records what remains.

\begin{lemma}\label{lem:face}
Suppose $\preceq$ satisfies strict (G), and a coherent model $\Dscr_K$ is
competitive on $S$. Then $K\subseteq\Delta_S\coloneqq\{m\in\Delta:m_j=0\text{ for }j\notin S\}$.
\end{lemma}
\begin{proof}
Suppose $m\in K$ has $m_j>0$ for some $j\notin S$. Consider gambles $g$ with
$g_s>0$ for every $s\in S$ and $g_j$ so negative that $m\cdot g<0$. They form a
non-null open set. $\Dscr_K$ does not accept them, since its lower prevision
is negative there. So on the corresponding problems its advice is silence or
acceptance of $-g$, and neither is good at any world of $S$.

Let $\Escr$ agree with $\Dscr_K$ except that it accepts exactly $g$ on these
problems. Then $\Escr$ is a strict guidance improvement of $\Dscr_K$ at every
world of $S$. By strict (G), $\Escr\prec_s\Dscr_K$ for all $s\in S$, so
$\Dscr_K$ is not competitive on $S$.
\end{proof}

\begin{corollary}\label{cor:face}
Under (G), (LA) and (S):
\begin{enumerate}[label=(\alph*)]
\item each vertex selection satisfies $C_{e_i}\approx D_i$, the ideal model;
\item each edge selection has credal set on that edge, so its advice depends
only on the two edge coordinates;
\item if the vacuous model $\Dscr_\Delta$ is undominated, it is tied at every
world only with selected models of full-support index.
\end{enumerate}
\end{corollary}
\begin{proof}
(a) and (b) follow from Lemma~\ref{lem:face} and v3's representation of
coherent sets. For (c), (KKMb) makes $\Dscr_\Delta$ competitive on $\supp p$
whenever it is tied with $C_p$. Lemma~\ref{lem:face} with $K=\Delta$ then
forces $\supp p$ to be all worlds.
\end{proof}

For a selected model $A$, let
$T_A\coloneqq\{p:C_p\text{ is tied with }A\text{ at every world}\}$.

\begin{lemma}\label{lem:fibres}
Assume (G), (LA) and (S).
\begin{enumerate}[label=(\alph*)]
\item $T_A=\bigcap_i\{p:C_p\preceq_iA\}$, and it is closed.
\item Along every path in $\Delta$ from $e_i$ to $e_j$, $i\ne j$, the selection
takes uncountably many values, no two tied at every world.
\item If every $\preceq_l$ is complete, then each selected $A$ has, at some
world $l$, no immediate strict predecessor: whenever $F\prec_lA$, some $G$
satisfies $F\prec_lG\prec_lA$.
\end{enumerate}
\end{lemma}
\begin{proof}
(a) If $C_p\preceq_iA$ at every world, undominatedness of $A$ makes each
comparison a tie. Closedness is part of (S).

(b) Let $\gamma:[0,1]\to\Delta$ be such a path. The preimages
$\gamma^{-1}(T_A)$, for distinct tie classes, are disjoint closed sets
covering $[0,1]$. At least two are nonempty: $C_{e_i}\approx D_i$ and
$C_{e_j}\approx D_j$ by Corollary~\ref{cor:face}, and $D_i\prec_iD_j$ by strict (G).
By Sierpi\'nski's theorem \cite{Sierpinski}, $[0,1]$ is not a countable union
of two or more pairwise disjoint nonempty closed sets. So there are
uncountably many tie classes.

(c) Suppose instead that at every world $l$ some $F_l\prec_lA$ has nothing
strictly between it and $A$. Take $q\in T_A$.
\begin{itemize}[leftmargin=*,itemsep=2pt]
\item Since $C_q\sim_lA\succ_lF_l$, the point $q$ lies outside the closed set
$\{p:C_p\preceq_lF_l\}$. So some neighbourhood $N$ of $q$ satisfies
$C_p\not\preceq_lF_l$ for every $l$ and $p\in N$.
\item Completeness gives $F_l\prec_lC_p$ for $p\in N$. Then $C_p\prec_lA$ is
impossible, since nothing lies strictly between $F_l$ and $A$, so
$A\preceq_lC_p$.
\item This holds at every world. Since $C_p$ is undominated, $C_p$ is tied
with $A$, so $N\subseteq T_A$.
\end{itemize}
Thus $T_A$ is open as well as closed. Since $\Delta$ is connected,
$T_A=\Delta$. This contradicts $D_1\prec_1D_2$.
\end{proof}

The next lemma generalizes the countability argument of BA's Proposition~6.
A \emph{primary} at world $l$ is a function $\Phi$ from models to $\R$ such
that $\Phi(\Dscr)<\Phi(\Escr)$ implies $\Dscr\prec_l\Escr$. A model $\Dscr$ has
an \emph{invisible improvement} at $l$ if some $\Escr\prec_l\Dscr$ has
$\Phi(\Escr)=\Phi(\Dscr)$.

\begin{lemma}\label{lem:primary}
Assume (S), and let $\Phi$ be a primary at world $l$.
\begin{enumerate}[label=(\alph*)]
\item The values $\Phi(C_p)$, over selected models with an invisible
improvement at $l$, form a countable set.
\item If $t<\Phi(C_p)$ and $t=\Phi(\Dscr)$ for some model $\Dscr$, then
$\Phi(C_q)\ge t$ for all $q$ near $p$.
\item If (LA) also holds, then undominated models with an invisible
improvement at $l$ take only countably many values of $\Phi$.
\end{enumerate}
\end{lemma}
\begin{proof}
(a) Let $\Escr$ be an invisible improvement of $C_p$. The point $p$ lies
outside the closed set $\{q:C_q\preceq_l\Escr\}$, so a neighbourhood $N$ of $p$
satisfies $C_q\not\preceq_l\Escr$. For $q\in N$ this forces
$\Phi(C_q)\ge\Phi(\Escr)=\Phi(C_p)$. So $p$ is a local minimizer of
$f(q)=\Phi(C_q)$.

The local-minimum values of a real function on a second-countable space form
a countable set. For each such value choose a basic open set $B$ on which it
is the minimum of $f$; distinct values receive distinct $B$, since the value
is $\inf_Bf$.

(b) Here $\Dscr\prec_lC_p$, so $p$ lies outside the closed set
$\{q:C_q\preceq_l\Dscr\}$. Near $p$, $C_q\not\preceq_l\Dscr$, which forces
$\Phi(C_q)\ge\Phi(\Dscr)$.

(c) Ties preserve $\Phi$: models with different $\Phi$-values are strictly
ordered. By (KKMb), an undominated model is tied with a selected model,
which then has the same $\Phi$-value and inherits the invisible improvement.
Apply (a).
\end{proof}

Part (c) is the general form of BA's countability argument. A ranking with a
real primary fails (S) as soon as uncountably many undominated coherent
models, with distinct primary values, admit invisible improvements.

\section{What remains}\label{sec:remaining}

Combine Theorems~\ref{thm:regular} and~\ref{thm:continuous} with
Section~\ref{sec:structure}. Suppose that, for some $n\ge3$, a ranking satisfies
(G), (IB), (LA) and (S) and leaves the vacuous model undominated. Then at some
world $i$ the following hold.
\begin{enumerate}[label=(\roman*),leftmargin=*]
\item Its comparisons between regular models have no additive weak extension.
In particular, on regular models it is not an intersection of complete,
independent, weakly monotone, continuous preorders.
\item If it is complete on regular models, it is discontinuous there. Every
selected model has no immediate strict predecessor at some world, and any
real primary obeys Lemma~\ref{lem:primary}.
\item If it is incomplete, every selected $C_p$ must still be comparable with
every model at some world of $\supp p$. BA's finite-sample theorem shows that
rankings comparing advice almost everywhere over finite samples of problems
are too sparse for this.
\end{enumerate}
Moreover, every vertex selects the ideal model, edge selections are supported
on their edge and fall into uncountably many tie classes along every edge,
and the vacuous model can be tied only with full-support selections.

Every construction we know falls into an excluded case:
\begin{itemize}[leftmargin=*,itemsep=2pt]
\item additive rankings (v3);
\item intersections of additive rankings, and complete continuous rankings of
any form (Theorem~\ref{thm:continuous});
\item lexicographic rankings with a real primary that is blind to some
improvements (Lemma~\ref{lem:primary}, BA);
\item sample-wise rankings (BA).
\end{itemize}
The benchmark rankings of v2 avoid all of these, but they compare models
through benchmark sets that depend on the whole model. For the ordinal version
this violates (IB) (v3, Proposition~1).

\paragraph{Conjecture.} For $n\ge3$, no ranking satisfies (G), (IB), (LA) and
(S) while leaving the vacuous model undominated.

A proof would have to handle incomplete rankings whose regular comparisons
are not intersections of complete continuous independent preorders.
Remark~\ref{rem:closed-graph} shows that such rankings exist, even with closed
graphs. A counterexample would have to be one of them, with comparisons dense
enough for competitiveness yet without a weak additive certificate.

\paragraph{Nearby reachable results.} Two complete results are already
available, each giving up one requirement.
\begin{itemize}[leftmargin=*,itemsep=2pt]
\item \emph{Without full coverage for $n\ge3$:} additive rankings satisfy every
requirement, including (IB) and (S). They cover all targets when $n=2$, and
the precise targets in every dimension (v3).
\item \emph{Without (IB):} the ordinal and real-valued constructions of v2
satisfy (E), (LA), (A), (S) and (C) with full coverage. They respect the
error-inclusion background order modulo null sets.
\end{itemize}
We know of no weakening of (IB) under which full coverage with (S) has been
proved.

\begin{thebibliography}{9}
\bibitem{Debreu} G.~Debreu, Topological methods in cardinal utility theory,
in K.~J. Arrow, S.~Karlin and P.~Suppes (eds.), \emph{Mathematical Methods in
the Social Sciences, 1959}, Stanford University Press, 1960, 16--26.
\bibitem{Sierpinski} W.~Sierpi\'nski, Un th\'eor\`eme sur les continus,
\emph{T\^ohoku Mathematical Journal} 13 (1918), 300--303.
\bibitem{Wakker} P.~P. Wakker, \emph{Additive Representations of Preferences},
Kluwer, Dordrecht, 1989.
\end{thebibliography}
\end{document}
